\scanpage{4}
\mspage{- 57 -}
n\textsuperscript{o}357
\[
f(x)\mathrel{\handwritten{\leq}}\handwritten{\sum}_{i=1}^{n}a_i\bigl[\varphi(s_i)+\handwritten{\eta}\bigr]g(s_i^{-1}x)
\]
et de même pour $f'$, et par suite :
\[
(f:g)+(f':g)\mathrel{\handwritten{\leq}}\handwritten{\sum}_{i=1}^{n}a_i\bigl[\varphi(s_i)+\varphi'(s_i)+2\handwritten{\eta}\bigr]\mathrel{\handwritten{\leq}}(1+2\handwritten{\eta})\handwritten{\sum}_{i=1}^{n}a_i
\]
puisque $\varphi+\varphi'\mathrel{\handwritten{\leq}}1$. En appliquant la définition de $[F:g]$, puis (ii) et (vii), on en conclut :
\[
(f:g)+(f':g)\mathrel{\handwritten{\leq}}(1+2\handwritten{\eta})(F:g)\mathrel{\handwritten{\leq}}(1+2\handwritten{\eta})\bigl[(f+f':g)+\handwritten{\delta}(h:g)\bigr]\ ;
\]
appliquant (v) à $f+f'$, $h$ et $g$, cela donne :
\[
(f:g)+(f':g)\mathrel{\handwritten{\leq}}(f+f':g)+2\handwritten{\eta}(f+f':h).(h:g)+(1+2\handwritten{\eta})\handwritten{\delta}.(h:g)\ .
\]
Si donc $f$, $f'$ et $h$ sont donnés comme dans (viii), il suffit de choisir $\handwritten{\delta}$ et $\handwritten{\eta}$ assez petits, puis de choisir $V$ comme il a été dit, pour qu'on ait l'inégalité formulée dans (viii).

Soit maintenant $\handwritten{\unsure{\mathfrak{F}|\mathscr{F}}}$ un ultrafiltre sur $\handwritten{\unsure{\mathcal{K}|\mathscr{K}}}^{\ast}_{+}$, plus fin que le filtre engendré par la réunion des $\handwritten{\mathcal{K}}^{\ast}_{+}(V)$ quand $V$ parcourt un système fondamental de voisinages compacts de $e$. D'autre part, fixons $f_{\unsure{o|0}}\mathrel{\handwritten{\in}}\handwritten{\mathcal{K}}^{\ast}_{+}$ ; posons, pour $f\mathrel{\handwritten{\in}}\handwritten{\mathcal{K}}^{\ast}_{+}$, $g\mathrel{\handwritten{\in}}\handwritten{\mathcal{K}}^{\ast}_{+}$ :
\[
I_g(f)=\frac{(f:g)}{(f_o:g)}\ ,
\]
puis $I(f)=\lim_{\handwritten{\mathfrak{F}}}I_g(f)$. D'après (vi), $I(f)$ est dans l'intervalle $[1/(f_o:f),(f:f_o)]$. D'après (ii), on a $I(cf)=c.I(f)$ pour $c>0$. D'après (vii), on a $I(f+f')\mathrel{\handwritten{\leq}}I(f)+I(f')$. D'après (viii), on a $I(f+f')\mathrel{\handwritten{\geq}}I(f)+I(f')-\varepsilon I(h)$ quel que soit $\varepsilon\mathrel{\handwritten{\geq}}0$ si $h$ est $\mathrel{\handwritten{\geq}}1$ sur le support de $f+f'$ ; comme on peut toujours choisir $h$ ainsi dans $\handwritten{\mathcal{K}}^{\ast}_{+}$, il s'ensuit que $I(f+f')=I(f)+I(f')$. Alors, d'après le Chap. II, \handwritten{\S}\ 2, n\textsuperscript{o}1, prop. 2, $I$ se prolonge (d'une manière et d'une seule) à une forme linéaire positive sur $\handwritten{\mathcal{K}}$, c'est-à-dire qu'il y a une mesure positive $\alpha$ sur $G$ telle que $I(f)=\handwritten{\int} f\,d\alpha$ pour $f\mathrel{\handwritten{\in}}\handwritten{\mathcal{K}}^{\ast}_{+}$. D'après (i), on a $I(\handwritten{\gamma}(\varepsilon_s)f)=I(f)$ quels que
% archive sheet number: 3

